\chapter{Entropion and Ectropion Repair}
\index{Entropion and Ectropion Repair}

\section*{Definition and Overview}
Entropion and ectropion are conditions of the eyelids that affect the way the eyelid sits against the eye. Entropion is an inward turning of the eyelid, so the lashes rub against the surface of the eye, causing irritation, tearing, and damage. Ectropion is an outward turning, or drooping, of the eyelid, leaving the inner surface exposed and causing dryness and watering. Both conditions most commonly affect the lower eyelid and are usually caused by age-related changes to the tissues, though injury, scarring, and other causes occur. Surgery is the definitive treatment.

\section*{Epidemiology}
Entropion and ectropion are common in older people, affecting a significant proportion of those over 60, and both are more common in men, particularly entropion. They are caused by the weakening of the muscles and tendons of the eyelid with age. Both conditions are often present for years before treatment, and they can cause significant eye irritation, tearing, and, if untreated, damage to the cornea.

\section*{Aetiology and Pathophysiology}
The eyelids are held in place by muscles, tendons, and elastic tissue, which weaken with age.
\begin{itemize}
    \item \textbf{Involutional (age-related) entropion:} the most common type, caused by laxity of the eyelid tissues and spasm of the eyelid muscle, turning the lid inward
    \item \textbf{Involutional ectropion:} age-related laxity makes the lid sag outward
    \item \textbf{Scarring:} injury, burns, or previous surgery can pull the lid inward or outward
    \item \textbf{Nerve damage:} facial nerve weakness, as in Bell's palsy, can cause ectropion
    \item \textbf{Effects of entropion:} the lashes and lid rub on the cornea and conjunctiva, causing abrasions and infection
    \item \textbf{Effects of ectropion:} the inner lid and eye are exposed, causing dryness, inflammation, and excessive tearing as the eye tries to compensate
\end{itemize}

\section*{Clinical Features}
The symptoms differ between the two conditions.
\begin{itemize}
    \item \textbf{Entropion:} a feeling of something in the eye, redness, tearing, pain, and sensitivity to light
    \item \textbf{Ectropion:} watering, dryness, redness, and a gritty sensation; the lid margin is visible, sagging away from the eye
    \item Tearing that runs down the cheek, common in both
    \item Crusting and discharge
    \item Blurred vision from the wet or damaged eye surface
    \item In entropion: visible inward turning of the lid and lashes
    \item In ectropion: visible drooping of the lower lid
\end{itemize}

\section*{Differential Diagnosis}
Other eye conditions cause similar symptoms.
\begin{itemize}
    \item Dry eye syndrome
    \item Blocked tear duct, causing watering
    \item Conjunctivitis and blepharitis
    \item Corneal abrasion and ulcer
    \item Trichiasis, in which lashes grow inwards without the lid turning
    \item Eyelid tumours
\end{itemize}

\section*{When to Seek Medical Care}
People with persistent eye irritation, tearing, or visible eyelid problems should be assessed by a GP and referred to an ophthalmologist. Surgery is usually recommended when the condition causes symptoms or threatens the cornea.

\textbf{Contact your local emergency services for:}
\begin{itemize}
    \item Sudden severe eye pain
    \item Sudden loss of vision
    \item A red eye with discharge and pain, which may indicate a serious infection
    \item Signs of corneal damage
\end{itemize}

\section*{Diagnosis and Investigations}
Diagnosis is made by examining the eye and eyelids.
\begin{itemize}
    \item \textbf{Examination:} assessment of the eyelid position, laxity, and muscle tone
    \item \textbf{Slit-lamp examination:} checking the cornea and eye surface for damage
    \item \textbf{Tear film assessment:} evaluating dryness or tearing
    \item \textbf{Assessment of the cause:} distinguishing age-related from scarring or nerve causes
\end{itemize}

\section*{Management}
Surgery corrects the eyelid position.
\begin{itemize}
    \item \textbf{Observation:} mild cases with few symptoms may be monitored
    \item \textbf{Lubricants:} drops and ointments protect the eye while awaiting surgery
    \item \textbf{Tape and Botox:} temporary measures for entropion in people unfit for surgery
    \item \textbf{Surgery for entropion:} tightening the eyelid and turning it outward, often a short procedure under local anaesthetic
    \item \textbf{Surgery for ectropion:} tightening the lax eyelid to return it to its normal position
    \item \textbf{Repair of the cause:} procedures to address scarring or nerve causes
    \item \textbf{Aftercare:} drops, ointment, and follow-up to monitor healing
\end{itemize}

\section*{Complications}
\begin{itemize}
    \item Corneal abrasions, ulcers, and scarring from untreated entropion
    \item Chronic conjunctivitis and dry eye from ectropion
    \item Infection
    \item Bleeding and bruising after surgery
    \item Over- or under-correction of the eyelid position
    \item Recurrence over time
\end{itemize}

\section*{Prognosis and Outcomes}
Surgery is very effective for both entropion and ectropion, and most people have a good cosmetic and functional result. The eye irritation and tearing resolve once the eyelid is correctly positioned. Complications are uncommon, and the procedures are usually performed as day surgery with a short recovery.

\section*{Prevention and Self-Care}
\begin{itemize}
    \item Use lubricating drops while waiting for surgery to protect the eye
    \item Protect the eye from injury and dust
    \item Follow the post-operative instructions, including drops and ointment
    \item Avoid rubbing the eyes
    \item Attend follow-up appointments
    \item Seek care promptly for any change in vision or increasing pain
\end{itemize}

\section*{Sources and Further Reading}
The following reputable sources provide further information for healthcare students and patients seeking reliable, current advice about entropion and ectropion.
\begin{itemize}
    \item American Academy of Ophthalmology. \textit{Entropion and ectropion patient information.} https://www.aao.org/eye-health/
    \item Healthdirect Australia. \textit{Eyelid problems.} https://www.healthdirect.gov.au/
    \item NHS. \textit{Ectropion and entropion.} https://www.nhs.uk/conditions/ectropion/
    \item Mayo Clinic. \textit{Ectropion: symptoms and treatment.} https://www.mayoclinic.org/diseases-conditions/ectropion/
    \item MedlinePlus. \textit{Entropion.} https://medlineplus.gov/ency/article/001009.htm
    \item The Royal Australian and New Zealand College of Ophthalmologists. \textit{Eyelid surgery information.} https://ranzco.edu/
\end{itemize}
